
\documentclass{article}
\usepackage{colortbl}
\usepackage{makecell}
\usepackage{multirow}
\usepackage{supertabular}

\begin{document}

\newcounter{utterance}

\centering \large Interaction Transcript for game `clean\_up', experiment `2\_hard\_7obj\_el', episode 1 with mistral{-}large{-}3{-}azure.
\vspace{24pt}

{ \footnotesize  \setcounter{utterance}{1}
\setlength{\tabcolsep}{0pt}
\begin{supertabular}{c@{$\;$}|p{.15\linewidth}@{}p{.15\linewidth}p{.15\linewidth}p{.15\linewidth}p{.15\linewidth}p{.15\linewidth}}
    \# & $\;$A & \multicolumn{4}{c}{Game Master} & $\;\:$B\\
    \hline

    \theutterance \stepcounter{utterance}  
    & & \multicolumn{4}{p{0.6\linewidth}}{
        \cellcolor[rgb]{0.9,0.9,0.9}{
            \makecell[{{p{\linewidth}}}]{
                \texttt{\tiny{[P1$\langle$GM]}}
                \texttt{Είμαι ο αρχηγός του παιχνιδιού σας και παίζετε ένα συνεργατικό παιχνίδι με το ακόλουθο πλέγμα ως πίνακα παιχνιδιού:} \\
\\ 
\\ 
\texttt{ 1234567} \\
\texttt{╔═══╤═╤═╗} \\
\texttt{║◌◌M│T│◌║ 1} \\
\texttt{║◌◌A└─┤S║ 2} \\
\texttt{║P◌◌O◌├─╢ 3} \\
\texttt{╟─┐◌┌─┤◌║ 4} \\
\texttt{║◌│D│◌│◌║ 5} \\
\texttt{╟─┼─┼─┼─╢ 6} \\
\texttt{║◌│◌│◌│◌║ 7} \\
\texttt{╚═╧═╧═╧═╝} \\
\\ 
\\ 
\\ 
\texttt{* Το άνω άκρο εμφανίζει συντεταγμένες x που αυξάνονται προς τα δεξιά και το δεξιό άκρο συντεταγμένες y που αυξάνονται προς τα κάτω.} \\
\texttt{* Τα ακόλουθα αντικείμενα τοποθετούνται τυχαία στο πλέγμα σας: 'P', 'O', 'T', 'S', 'D', 'A', 'M'.} \\
\\ 
\\ 
\texttt{Ο άλλος παίκτης βλέπει μια παραλλαγή του σκάφους του παιχνιδιού, όπου τα αντικείμενα τοποθετούνται σε διαφορετικές τυχαίες θέσεις. Δεν μπορείτε να δείτε το ταμπλό του άλλου παίκτη και εκείνος δεν μπορεί να δει το δικό σας.} \\
\\ 
\\ 
\texttt{**Στόχος:**} \\
\\ 
\texttt{Και οι δύο παίκτες πρέπει να μετακινήσουν τα αντικείμενα στο αντίστοιχο φόντο τους, έτσι ώστε πανομοιότυπα αντικείμενα να καταλήξουν στις ίδιες συντεταγμένες. Πρέπει να επικοινωνήσετε με τον άλλο παίκτη για να συμφωνήσετε σε μια κοινή διαμόρφωση του στόχου.} \\
\\ 
\\ 
\texttt{**Κανόνες:**} \\
\\ 
\texttt{* Σε κάθε γύρο, μπορείτε να στείλετε ακριβώς μία από τις ακόλουθες δύο εντολές:} \\
\texttt{ 1. `ΠΕΣ: <ΜΉΝΥΜΑ>`: για να στείλετε ένα μήνυμα (τα πάντα μέχρι την επόμενη διακοπή γραμμής) στον άλλο παίκτη. Θα το προωθήσω στον παρτενέρ σας.} \\
\texttt{2. `ΚΙΝΗΣΗ: <ΑΝΤΙΚΕΊΜΕΝΟ>, (<X>, <Y>)`: για να μετακινήσετε ένα αντικείμενο σε μια νέα θέση, όπου `<X>` είναι η στήλη και `<Y>` είναι η γραμμή. Θα σας ενημερώσω αν η μετακίνησή σας ήταν επιτυχής ή όχι.} \\
\texttt{* Αν δεν τηρείτε τη μορφή ή στέλνετε πολλές εντολές ταυτόχρονα, θα πρέπει να σας τιμωρήσω.} \\
\texttt{* Εάν και οι δύο παίκτες συγκεντρώσουν περισσότερες από 16 ποινές, χάνετε και οι δύο το παιχνίδι.} \\
\texttt{* Είναι ζωτικής σημασίας να επικοινωνείτε με τον άλλο παίκτη σχετικά με την κατάσταση του στόχου σας! Ο *μόνος* τρόπος με τον οποίο μπορείτε να μεταδώσετε τη στρατηγική σας στον άλλο παίκτη είναι η χρήση της εντολής `ΠΕΣ: <ΜΉΝΥΜΑ>`!} \\
\\ 
\\ 
\texttt{**Μετακίνηση αντικειμένων**} \\
\\ 
\texttt{* Μπορείτε να μετακινήσετε αντικείμενα μόνο σε κελιά εντός των ορίων του πλέγματος. Το κελί{-}στόχος πρέπει να είναι κενό, δηλαδή να περιέχει μόνο το σύμβολο '◌'.} \\
\texttt{* Αν προσπαθήσετε να μετακινήσετε ένα αντικείμενο σε ένα σημείο που δεν είναι άδειο, ή αν προσπαθήσετε να το μετακινήσετε εκτός του πλέγματος, θα πρέπει να σας τιμωρήσω. Μπορείτε να προσπαθήσετε ξανά.} \\
\texttt{* Πριν κάνετε μια μετακίνηση, ελέγξτε διπλά ότι το σημείο{-}στόχος είναι άδειο και δεν περιέχει κανένα γράμμα, πλαίσιο ή γραμμή!} \\
\\ 
\\ 
\texttt{**Τέλος του παιχνιδιού**} \\
\\ 
\texttt{Αν νομίζετε ότι έχετε πετύχει το στόχο της ευθυγράμμισης όλων των αντικειμένων, μπορείτε να ζητήσετε από τον άλλο παίκτη να τερματίσει το παιχνίδι στέλνοντας `ΠΕΣ: τελείωσες;`. Αν ο άλλος παίκτης σας ζητήσει να τερματίσετε το παιχνίδι και απαντήσετε με `ΠΕΣ: τελείωσα!`, το παιχνίδι θα τελειώσει.} \\
\\ 
\texttt{Και οι δύο παίκτες κερδίζουν αν το παιχνίδι τελειώσει μέσα σε 28 γύρους, όπου ένας γύρος ορίζεται ως δύο παίκτες που στέλνουν από μία έγκυρη εντολή.} \\
\\ 
\texttt{**Βαθμολογία:**} \\
\\ 
\texttt{Όσο πιο κοντά βρίσκονται τα ίδια αντικείμενα και στους δύο πίνακες του παιχνιδιού, τόσο περισσότερους πόντους παίρνετε. Οι ποινές μειώνουν τους πόντους σας. Μπορείτε να καταρρίψετε το ρεκόρ;} \\
\\ 
\\ 
\texttt{Παρακαλώ στείλτε ένα μήνυμα στον άλλο παίκτη για να ξεκινήσετε το παιχνίδι!} \\
            }
        }
    }
    & \\ \\

    \theutterance \stepcounter{utterance}  
    & \multicolumn{4}{p{0.6\linewidth}}{
        \cellcolor[rgb]{0.9,0.9,0.9}{
            \makecell[{{p{\linewidth}}}]{
                \texttt{\tiny{[P1$\rangle$GM]}}
                \texttt{Ας ξεκινήσουμε την επικοινωνία με τον άλλο παίκτη για να συμφωνήσουμε σε μια κοινή διαμόρφωση στόχου. Θα προτείνω μια στρατηγική για να τοποθετήσουμε τα αντικείμενα σε συγκεκριμένες συντεταγμένες, ώστε να είναι εύκολο να τα παρακολουθούμε και να τα ευθυγραμμίσουμε.} \\
\\ 
\texttt{**Πρώτο μήνυμα:**} \\
\texttt{```} \\
\texttt{ΠΕΣ: Γειά σου! Ας συμφωνήσουμε να τοποθετήσουμε τα αντικείμενα στις εξής θέσεις για να τα ευθυγραμμίσουμε:} \\
\texttt{{-} P: (1, 3)} \\
\texttt{{-} O: (4, 3)} \\
\texttt{{-} T: (4, 1)} \\
\texttt{{-} S: (6, 2)} \\
\texttt{{-} D: (3, 5)} \\
\texttt{{-} A: (3, 2)} \\
\texttt{{-} M: (3, 1)} \\
\\ 
\texttt{Εσύ έχεις κάποια πρόταση ή προτίμηση για συγκεκριμένες θέσεις;} \\
\texttt{```} \\
            }
        }
    }
    & & \\ \\

    \theutterance \stepcounter{utterance}  
    & & & \multicolumn{2}{p{0.3\linewidth}}{
        \cellcolor[rgb]{0.9,0.9,0.9}{
            \makecell[{{p{\linewidth}}}]{
                \texttt{\tiny{[GM$|$GM]}}
                \texttt{Ας ξεκινήσουμε την επικοινωνία με τον άλλο παίκτη για να συμφωνήσουμε σε μια κοινή διαμόρφωση στόχου. Θα προτείνω μια στρατηγική για να τοποθετήσουμε τα αντικείμενα σε συγκεκριμένες συντεταγμένες, ώστε να είναι εύκολο να τα παρακολουθούμε και να τα ευθυγραμμίσουμε.} \\
\\ 
\texttt{**Πρώτο μήνυμα:**} \\
\texttt{```} \\
\texttt{ΠΕΣ: Γειά σου! Ας συμφωνήσουμε να τοποθετήσουμε τα αντικείμενα στις εξής θέσεις για να τα ευθυγραμμίσουμε:} \\
\texttt{{-} P: (1, 3)} \\
\texttt{{-} O: (4, 3)} \\
\texttt{{-} T: (4, 1)} \\
\texttt{{-} S: (6, 2)} \\
\texttt{{-} D: (3, 5)} \\
\texttt{{-} A: (3, 2)} \\
\texttt{{-} M: (3, 1)} \\
\\ 
\texttt{Εσύ έχεις κάποια πρόταση ή προτίμηση για συγκεκριμένες θέσεις;} \\
\texttt{```} \\
            }
        }
    }
    & & \\ \\

    \theutterance \stepcounter{utterance}  
    & & & \multicolumn{2}{p{0.3\linewidth}}{
        \cellcolor[rgb]{0.9,0.9,0.9}{
            \makecell[{{p{\linewidth}}}]{
                \texttt{\tiny{[GM$|$GM]}}
                \texttt{Invalid format: head and tail are not empty} \\
\texttt{head: 'Ας ξεκινήσουμε την επικοινωνία με τον άλλο παίκτη για να συμφωνήσουμε σε μια κοινή διαμόρφωση στόχου. Θα προτείνω μια στρατηγική για να τοποθετήσουμε τα αντικείμενα σε συγκεκριμένες συντεταγμένες, ώστε να είναι εύκολο να τα παρακολουθούμε και να τα ευθυγραμμίσουμε.} \\
\\ 
\texttt{**Πρώτο μήνυμα:**} \\
\\ 
\texttt{'} \\
\texttt{tail: '} \\
\texttt{{-} P: (1, 3)} \\
\texttt{{-} O: (4, 3)} \\
\texttt{{-} T: (4, 1)} \\
\texttt{{-} S: (6, 2)} \\
\texttt{{-} D: (3, 5)} \\
\texttt{{-} A: (3, 2)} \\
\texttt{{-} M: (3, 1)} \\
\\ 
\texttt{Εσύ έχεις κάποια πρόταση ή προτίμηση για συγκεκριμένες θέσεις;'} \\
            }
        }
    }
    & & \\ \\

    \theutterance \stepcounter{utterance}  
    & & \multicolumn{4}{p{0.6\linewidth}}{
        \cellcolor[rgb]{0.9,0.9,0.9}{
            \makecell[{{p{\linewidth}}}]{
                \texttt{\tiny{[P1$\langle$GM]}}
                \texttt{Ποινή: Το μήνυμά σας δεν πρέπει να περιέχει τίποτα πριν ή μετά την εντολή!} \\
\\ 
\texttt{Βεβαιωθείτε ότι η απάντησή σας περιέχει μόνο `ΠΕΣ: <ΜΉΝΥΜΑ>` ή `ΚΙΝΗΣΗ: <ΑΝΤΙΚΕΊΜΕΝΟ>, (<X>, <Y>)`, και τίποτα άλλο!} \\
\\ 
\texttt{Έχετε συσσωρεύσει συλλογικά 1 από 16 ποινές.} \\
\texttt{Παρακαλώ δοκιμάστε ξανά!} \\
            }
        }
    }
    & \\ \\

    \theutterance \stepcounter{utterance}  
    & \multicolumn{4}{p{0.6\linewidth}}{
        \cellcolor[rgb]{0.9,0.9,0.9}{
            \makecell[{{p{\linewidth}}}]{
                \texttt{\tiny{[P1$\rangle$GM]}}
                \texttt{```} \\
\texttt{ΠΕΣ: Γειά σου! Ας βάλουμε τα αντικείμενα στις εξής θέσεις:} \\
\texttt{P: (1,3), O: (4,3), T: (4,1), S: (6,2), D: (3,5), A: (3,2), M: (3,1).} \\
\texttt{Συμφωνείς;} \\
\texttt{```} \\
            }
        }
    }
    & & \\ \\

    \theutterance \stepcounter{utterance}  
    & & & \multicolumn{2}{p{0.3\linewidth}}{
        \cellcolor[rgb]{0.9,0.9,0.9}{
            \makecell[{{p{\linewidth}}}]{
                \texttt{\tiny{[GM$|$GM]}}
                \texttt{```} \\
\texttt{ΠΕΣ: Γειά σου! Ας βάλουμε τα αντικείμενα στις εξής θέσεις:} \\
\texttt{P: (1,3), O: (4,3), T: (4,1), S: (6,2), D: (3,5), A: (3,2), M: (3,1).} \\
\texttt{Συμφωνείς;} \\
\texttt{```} \\
            }
        }
    }
    & & \\ \\

    \theutterance \stepcounter{utterance}  
    & & & \multicolumn{2}{p{0.3\linewidth}}{
        \cellcolor[rgb]{0.9,0.9,0.9}{
            \makecell[{{p{\linewidth}}}]{
                \texttt{\tiny{[GM$|$GM]}}
                \texttt{Invalid format: tail is not empty: '} \\
\texttt{P: (1,3), O: (4,3), T: (4,1), S: (6,2), D: (3,5), A: (3,2), M: (3,1).} \\
\texttt{Συμφωνείς;} \\
            }
        }
    }
    & & \\ \\

    \theutterance \stepcounter{utterance}  
    & & \multicolumn{4}{p{0.6\linewidth}}{
        \cellcolor[rgb]{0.9,0.9,0.9}{
            \makecell[{{p{\linewidth}}}]{
                \texttt{\tiny{[P1$\langle$GM]}}
                \texttt{Ποινή: Το μήνυμά σας δεν πρέπει να περιέχει τίποτα μετά την εντολή!} \\
\\ 
\texttt{Βεβαιωθείτε ότι η απάντησή σας περιέχει μόνο `ΠΕΣ: <ΜΉΝΥΜΑ>` ή `ΚΙΝΗΣΗ: <ΑΝΤΙΚΕΊΜΕΝΟ>, (<X>, <Y>)`, και τίποτα άλλο!} \\
\\ 
\texttt{Έχετε συσσωρεύσει συλλογικά 2 από 16 ποινές.} \\
\texttt{Παρακαλώ δοκιμάστε ξανά!} \\
            }
        }
    }
    & \\ \\

    \theutterance \stepcounter{utterance}  
    & \multicolumn{4}{p{0.6\linewidth}}{
        \cellcolor[rgb]{0.9,0.9,0.9}{
            \makecell[{{p{\linewidth}}}]{
                \texttt{\tiny{[P1$\rangle$GM]}}
                \texttt{```} \\
\texttt{ΠΕΣ: Γειά σου! Προτείνω P(1,3), O(4,3), T(4,1), S(6,2), D(3,5), A(3,2), M(3,1). Συμφωνείς;} \\
\texttt{```} \\
            }
        }
    }
    & & \\ \\

    \theutterance \stepcounter{utterance}  
    & & & \multicolumn{2}{p{0.3\linewidth}}{
        \cellcolor[rgb]{0.9,0.9,0.9}{
            \makecell[{{p{\linewidth}}}]{
                \texttt{\tiny{[GM$|$GM]}}
                \texttt{```} \\
\texttt{ΠΕΣ: Γειά σου! Προτείνω P(1,3), O(4,3), T(4,1), S(6,2), D(3,5), A(3,2), M(3,1). Συμφωνείς;} \\
\texttt{```} \\
            }
        }
    }
    & & \\ \\

    \theutterance \stepcounter{utterance}  
    & & \multicolumn{4}{p{0.6\linewidth}}{
        \cellcolor[rgb]{0.9,0.9,0.9}{
            \makecell[{{p{\linewidth}}}]{
                \texttt{\tiny{[GM$\rangle$P2]}}
                \texttt{Είμαι ο αρχηγός του παιχνιδιού σας και παίζετε ένα συνεργατικό παιχνίδι με το ακόλουθο πλέγμα ως πίνακα παιχνιδιού:} \\
\\ 
\\ 
\texttt{ 1234567} \\
\texttt{╔═══╤═╤═╗} \\
\texttt{║◌◌◌│◌│◌║ 1} \\
\texttt{║T◌M└─┤◌║ 2} \\
\texttt{║AS◌P◌├─╢ 3} \\
\texttt{╟─┐D┌─┤O║ 4} \\
\texttt{║◌│◌│◌│◌║ 5} \\
\texttt{╟─┼─┼─┼─╢ 6} \\
\texttt{║◌│◌│◌│◌║ 7} \\
\texttt{╚═╧═╧═╧═╝} \\
\\ 
\\ 
\\ 
\texttt{* Το άνω άκρο εμφανίζει συντεταγμένες x που αυξάνονται προς τα δεξιά και το δεξιό άκρο συντεταγμένες y που αυξάνονται προς τα κάτω.} \\
\texttt{* Τα ακόλουθα αντικείμενα τοποθετούνται τυχαία στο πλέγμα σας: 'P', 'O', 'T', 'S', 'D', 'A', 'M'.} \\
\\ 
\\ 
\texttt{Ο άλλος παίκτης βλέπει μια παραλλαγή του σκάφους του παιχνιδιού, όπου τα αντικείμενα τοποθετούνται σε διαφορετικές τυχαίες θέσεις. Δεν μπορείτε να δείτε το ταμπλό του άλλου παίκτη και εκείνος δεν μπορεί να δει το δικό σας.} \\
\\ 
\\ 
\texttt{**Στόχος:**} \\
\\ 
\texttt{Και οι δύο παίκτες πρέπει να μετακινήσουν τα αντικείμενα στο αντίστοιχο φόντο τους, έτσι ώστε πανομοιότυπα αντικείμενα να καταλήξουν στις ίδιες συντεταγμένες. Πρέπει να επικοινωνήσετε με τον άλλο παίκτη για να συμφωνήσετε σε μια κοινή διαμόρφωση του στόχου.} \\
\\ 
\\ 
\texttt{**Κανόνες:**} \\
\\ 
\texttt{* Σε κάθε γύρο, μπορείτε να στείλετε ακριβώς μία από τις ακόλουθες δύο εντολές:} \\
\texttt{ 1. `ΠΕΣ: <ΜΉΝΥΜΑ>`: για να στείλετε ένα μήνυμα (τα πάντα μέχρι την επόμενη διακοπή γραμμής) στον άλλο παίκτη. Θα το προωθήσω στον παρτενέρ σας.} \\
\texttt{2. `ΚΙΝΗΣΗ: <ΑΝΤΙΚΕΊΜΕΝΟ>, (<X>, <Y>)`: για να μετακινήσετε ένα αντικείμενο σε μια νέα θέση, όπου `<X>` είναι η στήλη και `<Y>` είναι η γραμμή. Θα σας ενημερώσω αν η μετακίνησή σας ήταν επιτυχής ή όχι.} \\
\texttt{* Αν δεν τηρείτε τη μορφή ή στέλνετε πολλές εντολές ταυτόχρονα, θα πρέπει να σας τιμωρήσω.} \\
\texttt{* Εάν και οι δύο παίκτες συγκεντρώσουν περισσότερες από 16 ποινές, χάνετε και οι δύο το παιχνίδι.} \\
\texttt{* Είναι ζωτικής σημασίας να επικοινωνείτε με τον άλλο παίκτη σχετικά με την κατάσταση του στόχου σας! Ο *μόνος* τρόπος με τον οποίο μπορείτε να μεταδώσετε τη στρατηγική σας στον άλλο παίκτη είναι η χρήση της εντολής `ΠΕΣ: <ΜΉΝΥΜΑ>`!} \\
\\ 
\\ 
\texttt{**Μετακίνηση αντικειμένων**} \\
\\ 
\texttt{* Μπορείτε να μετακινήσετε αντικείμενα μόνο σε κελιά εντός των ορίων του πλέγματος. Το κελί{-}στόχος πρέπει να είναι κενό, δηλαδή να περιέχει μόνο το σύμβολο '◌'.} \\
\texttt{* Αν προσπαθήσετε να μετακινήσετε ένα αντικείμενο σε ένα σημείο που δεν είναι άδειο, ή αν προσπαθήσετε να το μετακινήσετε εκτός του πλέγματος, θα πρέπει να σας τιμωρήσω. Μπορείτε να προσπαθήσετε ξανά.} \\
\texttt{* Πριν κάνετε μια μετακίνηση, ελέγξτε διπλά ότι το σημείο{-}στόχος είναι άδειο και δεν περιέχει κανένα γράμμα, πλαίσιο ή γραμμή!} \\
\\ 
\\ 
\texttt{**Τέλος του παιχνιδιού**} \\
\\ 
\texttt{Αν νομίζετε ότι έχετε πετύχει το στόχο της ευθυγράμμισης όλων των αντικειμένων, μπορείτε να ζητήσετε από τον άλλο παίκτη να τερματίσει το παιχνίδι στέλνοντας `ΠΕΣ: τελείωσες;`. Αν ο άλλος παίκτης σας ζητήσει να τερματίσετε το παιχνίδι και απαντήσετε με `ΠΕΣ: τελείωσα!`, το παιχνίδι θα τελειώσει.} \\
\\ 
\texttt{Και οι δύο παίκτες κερδίζουν αν το παιχνίδι τελειώσει μέσα σε 28 γύρους, όπου ένας γύρος ορίζεται ως δύο παίκτες που στέλνουν από μία έγκυρη εντολή.} \\
\\ 
\texttt{**Βαθμολογία:**} \\
\\ 
\texttt{Όσο πιο κοντά βρίσκονται τα ίδια αντικείμενα και στους δύο πίνακες του παιχνιδιού, τόσο περισσότερους πόντους παίρνετε. Οι ποινές μειώνουν τους πόντους σας. Μπορείτε να καταρρίψετε το ρεκόρ;} \\
\\ 
\texttt{Ο άλλος παίκτης ξεκίνησε το παιχνίδι στέλνοντας αυτό το μήνυμα:} \\
\\ 
\texttt{"Γειά σου! Προτείνω P(1,3), O(4,3), T(4,1), S(6,2), D(3,5), A(3,2), M(3,1). Συμφωνείς;"} \\
\\ 
\texttt{Ποια είναι η πρώτη σας εντολή;} \\
            }
        }
    }
    & \\ \\

    \theutterance \stepcounter{utterance}  
    & & & \multicolumn{4}{p{0.6\linewidth}}{
        \cellcolor[rgb]{0.9,0.9,0.9}{
            \makecell[{{p{\linewidth}}}]{
                \texttt{\tiny{[GM$\langle$P2]}}
                \texttt{Απαντώντας στο μήνυμα του άλλου παίκτη, η πρώτη μου εντολή θα είναι να **συμφωνήσω με την πρότασή του**, καθώς φαίνεται λογική και καλύπτει όλες τις συντεταγμένες χωρίς συγκρούσεις (σύμφωνα με το δικό μου ταμπλό). Θα στείλω:} \\
\\ 
\texttt{```} \\
\texttt{ΠΕΣ: Ναι, συμφωνώ! Ας ξεκινήσουμε με τις θέσεις: P(1,3), O(4,3), T(4,1), S(6,2), D(3,5), A(3,2), M(3,1). Θα μετακινήσω πρώτα το 'A' στη θέση του.} \\
\texttt{```} \\
\\ 
\texttt{\#\#\# Επόμενο βήμα:} \\
\texttt{Μετά την αποστολή του μηνύματος, θα προχωρήσω στην **πρώτη μετακίνηση** του `'A'` στη συμφωνημένη θέση `(3,2)`. Η εντολή θα είναι:} \\
\texttt{```} \\
\texttt{ΚΙΝΗΣΗ: A, (3, 2)} \\
\texttt{```} \\
\\ 
\texttt{\#\#\# Σημειώσεις:} \\
\texttt{1. **Επαλήθευση κενού κελιού**: Στο δικό μου ταμπλό, το `(3,2)` είναι άδειο (`◌`), οπότε η μετακίνηση είναι έγκυρη.} \\
\texttt{2. **Συντονισμός**: Αν ο άλλος παίκτης έχει διαφορετική αρχική θέση για το `'A'`, θα πρέπει να το μετακινήσει και αυτός στο `(3,2)` στους επόμενους γύρους.} \\
\texttt{3. **Στρατηγική**: Θα εστιάσω σε αντικείμενα που είναι κοντά στις συμφωνημένες θέσεις για να ελαχιστοποιήσω τις κινήσεις.} \\
\\ 
\texttt{Αν η μετακίνηση είναι επιτυχής, θα συνεχίσω με το επόμενο αντικείμενο (π.χ. `'M'` ή `'D'`). Αν όχι, θα ζητήσω διευκρινίσεις.} \\
            }
        }
    }
     \\ \\

    \theutterance \stepcounter{utterance}  
    & & & \multicolumn{2}{p{0.3\linewidth}}{
        \cellcolor[rgb]{0.9,0.9,0.9}{
            \makecell[{{p{\linewidth}}}]{
                \texttt{\tiny{[GM$|$GM]}}
                \texttt{Απαντώντας στο μήνυμα του άλλου παίκτη, η πρώτη μου εντολή θα είναι να **συμφωνήσω με την πρότασή του**, καθώς φαίνεται λογική και καλύπτει όλες τις συντεταγμένες χωρίς συγκρούσεις (σύμφωνα με το δικό μου ταμπλό). Θα στείλω:} \\
\\ 
\texttt{```} \\
\texttt{ΠΕΣ: Ναι, συμφωνώ! Ας ξεκινήσουμε με τις θέσεις: P(1,3), O(4,3), T(4,1), S(6,2), D(3,5), A(3,2), M(3,1). Θα μετακινήσω πρώτα το 'A' στη θέση του.} \\
\texttt{```} \\
\\ 
\texttt{\#\#\# Επόμενο βήμα:} \\
\texttt{Μετά την αποστολή του μηνύματος, θα προχωρήσω στην **πρώτη μετακίνηση** του `'A'` στη συμφωνημένη θέση `(3,2)`. Η εντολή θα είναι:} \\
\texttt{```} \\
\texttt{ΚΙΝΗΣΗ: A, (3, 2)} \\
\texttt{```} \\
\\ 
\texttt{\#\#\# Σημειώσεις:} \\
\texttt{1. **Επαλήθευση κενού κελιού**: Στο δικό μου ταμπλό, το `(3,2)` είναι άδειο (`◌`), οπότε η μετακίνηση είναι έγκυρη.} \\
\texttt{2. **Συντονισμός**: Αν ο άλλος παίκτης έχει διαφορετική αρχική θέση για το `'A'`, θα πρέπει να το μετακινήσει και αυτός στο `(3,2)` στους επόμενους γύρους.} \\
\texttt{3. **Στρατηγική**: Θα εστιάσω σε αντικείμενα που είναι κοντά στις συμφωνημένες θέσεις για να ελαχιστοποιήσω τις κινήσεις.} \\
\\ 
\texttt{Αν η μετακίνηση είναι επιτυχής, θα συνεχίσω με το επόμενο αντικείμενο (π.χ. `'M'` ή `'D'`). Αν όχι, θα ζητήσω διευκρινίσεις.} \\
            }
        }
    }
    & & \\ \\

    \theutterance \stepcounter{utterance}  
    & & & \multicolumn{2}{p{0.3\linewidth}}{
        \cellcolor[rgb]{0.9,0.9,0.9}{
            \makecell[{{p{\linewidth}}}]{
                \texttt{\tiny{[GM$|$GM]}}
                \texttt{Invalid response format: Απαντώντας στο μήνυμα του άλλου παίκτη, η πρώτη μου εντολή θα είναι να **συμφωνήσω με την πρότασή του**, καθώς φαίνεται λογική και καλύπτει όλες τις συντεταγμένες χωρίς συγκρούσεις (σύμφωνα με το δικό μου ταμπλό). Θα στείλω:} \\
\\ 
\\ 
\texttt{ΠΕΣ: Ναι, συμφωνώ! Ας ξεκινήσουμε με τις θέσεις: P(1,3), O(4,3), T(4,1), S(6,2), D(3,5), A(3,2), M(3,1). Θα μετακινήσω πρώτα το 'A' στη θέση του.} \\
\\ 
\\ 
\texttt{\#\#\# Επόμενο βήμα:} \\
\texttt{Μετά την αποστολή του μηνύματος, θα προχωρήσω στην **πρώτη μετακίνηση** του 'A' στη συμφωνημένη θέση (3,2). Η εντολή θα είναι:} \\
\\ 
\texttt{ΚΙΝΗΣΗ: A, (3, 2)} \\
\\ 
\\ 
\texttt{\#\#\# Σημειώσεις:} \\
\texttt{1. **Επαλήθευση κενού κελιού**: Στο δικό μου ταμπλό, το (3,2) είναι άδειο (◌), οπότε η μετακίνηση είναι έγκυρη.} \\
\texttt{2. **Συντονισμός**: Αν ο άλλος παίκτης έχει διαφορετική αρχική θέση για το 'A', θα πρέπει να το μετακινήσει και αυτός στο (3,2) στους επόμενους γύρους.} \\
\texttt{3. **Στρατηγική**: Θα εστιάσω σε αντικείμενα που είναι κοντά στις συμφωνημένες θέσεις για να ελαχιστοποιήσω τις κινήσεις.} \\
\\ 
\texttt{Αν η μετακίνηση είναι επιτυχής, θα συνεχίσω με το επόμενο αντικείμενο (π.χ. 'M' ή 'D'). Αν όχι, θα ζητήσω διευκρινίσεις.} \\
            }
        }
    }
    & & \\ \\

    \theutterance \stepcounter{utterance}  
    & & \multicolumn{4}{p{0.6\linewidth}}{
        \cellcolor[rgb]{0.9,0.9,0.9}{
            \makecell[{{p{\linewidth}}}]{
                \texttt{\tiny{[GM$\rangle$P2]}}
                \texttt{Ποινή: Το μήνυμά σας περιέχει περισσότερες από μία εντολές!} \\
\\ 
\texttt{Βεβαιωθείτε ότι η απάντησή σας περιέχει μόνο `ΠΕΣ: <ΜΉΝΥΜΑ>` ή `ΚΙΝΗΣΗ: <ΑΝΤΙΚΕΊΜΕΝΟ>, (<X>, <Y>)`, και τίποτα άλλο!} \\
\\ 
\texttt{Έχετε συσσωρεύσει συλλογικά 3 από 16 ποινές.} \\
\texttt{Παρακαλώ δοκιμάστε ξανά!} \\
            }
        }
    }
    & \\ \\

    \theutterance \stepcounter{utterance}  
    & & & \multicolumn{4}{p{0.6\linewidth}}{
        \cellcolor[rgb]{0.9,0.9,0.9}{
            \makecell[{{p{\linewidth}}}]{
                \texttt{\tiny{[GM$\langle$P2]}}
                \texttt{Συνεχίζω με **μόνο μία εντολή**, όπως απαιτείται. Επιλέγω να απαντήσω στο μήνυμα του άλλου παίκτη:} \\
\\ 
\texttt{```} \\
\texttt{ΠΕΣ: Ναι, συμφωνώ με τις θέσεις!} \\
\texttt{```} \\
\\ 
\texttt{Αυτή είναι η πρώτη μου έγκυρη κίνηση. Θα περιμένω την επιβεβαίωση του άλλου παίκτη πριν προχωρήσω σε μετακίνηση αντικειμένου.} \\
            }
        }
    }
     \\ \\

    \theutterance \stepcounter{utterance}  
    & & & \multicolumn{2}{p{0.3\linewidth}}{
        \cellcolor[rgb]{0.9,0.9,0.9}{
            \makecell[{{p{\linewidth}}}]{
                \texttt{\tiny{[GM$|$GM]}}
                \texttt{Συνεχίζω με **μόνο μία εντολή**, όπως απαιτείται. Επιλέγω να απαντήσω στο μήνυμα του άλλου παίκτη:} \\
\\ 
\texttt{```} \\
\texttt{ΠΕΣ: Ναι, συμφωνώ με τις θέσεις!} \\
\texttt{```} \\
\\ 
\texttt{Αυτή είναι η πρώτη μου έγκυρη κίνηση. Θα περιμένω την επιβεβαίωση του άλλου παίκτη πριν προχωρήσω σε μετακίνηση αντικειμένου.} \\
            }
        }
    }
    & & \\ \\

    \theutterance \stepcounter{utterance}  
    & & & \multicolumn{2}{p{0.3\linewidth}}{
        \cellcolor[rgb]{0.9,0.9,0.9}{
            \makecell[{{p{\linewidth}}}]{
                \texttt{\tiny{[GM$|$GM]}}
                \texttt{Invalid format: head and tail are not empty} \\
\texttt{head: 'Συνεχίζω με **μόνο μία εντολή**, όπως απαιτείται. Επιλέγω να απαντήσω στο μήνυμα του άλλου παίκτη:} \\
\\ 
\\ 
\texttt{'} \\
\texttt{tail: '} \\
\\ 
\\ 
\texttt{Αυτή είναι η πρώτη μου έγκυρη κίνηση. Θα περιμένω την επιβεβαίωση του άλλου παίκτη πριν προχωρήσω σε μετακίνηση αντικειμένου.'} \\
            }
        }
    }
    & & \\ \\

    \theutterance \stepcounter{utterance}  
    & & \multicolumn{4}{p{0.6\linewidth}}{
        \cellcolor[rgb]{0.9,0.9,0.9}{
            \makecell[{{p{\linewidth}}}]{
                \texttt{\tiny{[GM$\rangle$P2]}}
                \texttt{Ποινή: Το μήνυμά σας δεν πρέπει να περιέχει τίποτα πριν ή μετά την εντολή!} \\
\\ 
\texttt{Βεβαιωθείτε ότι η απάντησή σας περιέχει μόνο `ΠΕΣ: <ΜΉΝΥΜΑ>` ή `ΚΙΝΗΣΗ: <ΑΝΤΙΚΕΊΜΕΝΟ>, (<X>, <Y>)`, και τίποτα άλλο!} \\
\\ 
\texttt{Έχετε συσσωρεύσει συλλογικά 4 από 16 ποινές.} \\
\texttt{Παρακαλώ δοκιμάστε ξανά!} \\
            }
        }
    }
    & \\ \\

    \theutterance \stepcounter{utterance}  
    & & & \multicolumn{4}{p{0.6\linewidth}}{
        \cellcolor[rgb]{0.9,0.9,0.9}{
            \makecell[{{p{\linewidth}}}]{
                \texttt{\tiny{[GM$\langle$P2]}}
                \texttt{Σωστά, καταλαβαίνω. Η επόμενη εντολή μου είναι:} \\
\\ 
\texttt{```} \\
\texttt{ΠΕΣ: Συμφωνώ με τις θέσεις P(1,3), O(4,3), T(4,1), S(6,2), D(3,5), A(3,2), M(3,1).} \\
\texttt{```} \\
            }
        }
    }
     \\ \\

    \theutterance \stepcounter{utterance}  
    & & & \multicolumn{2}{p{0.3\linewidth}}{
        \cellcolor[rgb]{0.9,0.9,0.9}{
            \makecell[{{p{\linewidth}}}]{
                \texttt{\tiny{[GM$|$GM]}}
                \texttt{Σωστά, καταλαβαίνω. Η επόμενη εντολή μου είναι:} \\
\\ 
\texttt{```} \\
\texttt{ΠΕΣ: Συμφωνώ με τις θέσεις P(1,3), O(4,3), T(4,1), S(6,2), D(3,5), A(3,2), M(3,1).} \\
\texttt{```} \\
            }
        }
    }
    & & \\ \\

    \theutterance \stepcounter{utterance}  
    & & & \multicolumn{2}{p{0.3\linewidth}}{
        \cellcolor[rgb]{0.9,0.9,0.9}{
            \makecell[{{p{\linewidth}}}]{
                \texttt{\tiny{[GM$|$GM]}}
                \texttt{Invalid format: head is not empty: 'Σωστά, καταλαβαίνω. Η επόμενη εντολή μου είναι:} \\
\\ 
\\ 
\texttt{'} \\
            }
        }
    }
    & & \\ \\

    \theutterance \stepcounter{utterance}  
    & & \multicolumn{4}{p{0.6\linewidth}}{
        \cellcolor[rgb]{0.9,0.9,0.9}{
            \makecell[{{p{\linewidth}}}]{
                \texttt{\tiny{[GM$\rangle$P2]}}
                \texttt{Ποινή: Το μήνυμά σας δεν πρέπει να περιέχει τίποτα πριν από την εντολή!} \\
\\ 
\texttt{Βεβαιωθείτε ότι η απάντησή σας περιέχει μόνο `ΠΕΣ: <ΜΉΝΥΜΑ>` ή `ΚΙΝΗΣΗ: <ΑΝΤΙΚΕΊΜΕΝΟ>, (<X>, <Y>)`, και τίποτα άλλο!} \\
\\ 
\texttt{Έχετε συσσωρεύσει συλλογικά 5 από 16 ποινές.} \\
\texttt{Παρακαλώ δοκιμάστε ξανά!} \\
            }
        }
    }
    & \\ \\

    \theutterance \stepcounter{utterance}  
    & & & \multicolumn{4}{p{0.6\linewidth}}{
        \cellcolor[rgb]{0.9,0.9,0.9}{
            \makecell[{{p{\linewidth}}}]{
                \texttt{\tiny{[GM$\langle$P2]}}
                \texttt{Συγγνώμη για τη σύγχυση. Η σωστή εντολή είναι:} \\
\\ 
\texttt{```} \\
\texttt{ΠΕΣ: Συμφωνώ} \\
\texttt{```} \\
            }
        }
    }
     \\ \\

    \theutterance \stepcounter{utterance}  
    & & & \multicolumn{2}{p{0.3\linewidth}}{
        \cellcolor[rgb]{0.9,0.9,0.9}{
            \makecell[{{p{\linewidth}}}]{
                \texttt{\tiny{[GM$|$GM]}}
                \texttt{Συγγνώμη για τη σύγχυση. Η σωστή εντολή είναι:} \\
\\ 
\texttt{```} \\
\texttt{ΠΕΣ: Συμφωνώ} \\
\texttt{```} \\
            }
        }
    }
    & & \\ \\

    \theutterance \stepcounter{utterance}  
    & & & \multicolumn{2}{p{0.3\linewidth}}{
        \cellcolor[rgb]{0.9,0.9,0.9}{
            \makecell[{{p{\linewidth}}}]{
                \texttt{\tiny{[GM$|$GM]}}
                \texttt{Invalid format: head is not empty: 'Συγγνώμη για τη σύγχυση. Η σωστή εντολή είναι:} \\
\\ 
\\ 
\texttt{'} \\
            }
        }
    }
    & & \\ \\

    \theutterance \stepcounter{utterance}  
    & & \multicolumn{4}{p{0.6\linewidth}}{
        \cellcolor[rgb]{0.9,0.9,0.9}{
            \makecell[{{p{\linewidth}}}]{
                \texttt{\tiny{[GM$\rangle$P2]}}
                \texttt{Ποινή: Το μήνυμά σας δεν πρέπει να περιέχει τίποτα πριν από την εντολή!} \\
\\ 
\texttt{Βεβαιωθείτε ότι η απάντησή σας περιέχει μόνο `ΠΕΣ: <ΜΉΝΥΜΑ>` ή `ΚΙΝΗΣΗ: <ΑΝΤΙΚΕΊΜΕΝΟ>, (<X>, <Y>)`, και τίποτα άλλο!} \\
\\ 
\texttt{Έχετε συσσωρεύσει συλλογικά 6 από 16 ποινές.} \\
\texttt{Παρακαλώ δοκιμάστε ξανά!} \\
            }
        }
    }
    & \\ \\

    \theutterance \stepcounter{utterance}  
    & & & \multicolumn{4}{p{0.6\linewidth}}{
        \cellcolor[rgb]{0.9,0.9,0.9}{
            \makecell[{{p{\linewidth}}}]{
                \texttt{\tiny{[GM$\langle$P2]}}
                \texttt{Δοκιμάζω ξανά με **αυστηρά μόνο την εντολή**:} \\
\\ 
\texttt{```} \\
\texttt{ΠΕΣ: ΟΚ για τις θέσεις.} \\
\texttt{```} \\
\\ 
\texttt{*(Σημείωση: Αν αυτό προκαλέσει νέα ποινή, θα στείλω απλά `ΠΕΣ: ΟΚ` στον επόμενο γύρο.)*} \\
            }
        }
    }
     \\ \\

    \theutterance \stepcounter{utterance}  
    & & & \multicolumn{2}{p{0.3\linewidth}}{
        \cellcolor[rgb]{0.9,0.9,0.9}{
            \makecell[{{p{\linewidth}}}]{
                \texttt{\tiny{[GM$|$GM]}}
                \texttt{Δοκιμάζω ξανά με **αυστηρά μόνο την εντολή**:} \\
\\ 
\texttt{```} \\
\texttt{ΠΕΣ: ΟΚ για τις θέσεις.} \\
\texttt{```} \\
\\ 
\texttt{*(Σημείωση: Αν αυτό προκαλέσει νέα ποινή, θα στείλω απλά `ΠΕΣ: ΟΚ` στον επόμενο γύρο.)*} \\
            }
        }
    }
    & & \\ \\

    \theutterance \stepcounter{utterance}  
    & & & \multicolumn{2}{p{0.3\linewidth}}{
        \cellcolor[rgb]{0.9,0.9,0.9}{
            \makecell[{{p{\linewidth}}}]{
                \texttt{\tiny{[GM$|$GM]}}
                \texttt{Invalid format: head is not empty: 'Δοκιμάζω ξανά με **αυστηρά μόνο την εντολή**:} \\
\\ 
\\ 
\texttt{ΠΕΣ: ΟΚ για τις θέσεις.} \\
\\ 
\\ 
\texttt{*(Σημείωση: Αν αυτό προκαλέσει νέα ποινή, θα στείλω απλά '} \\
            }
        }
    }
    & & \\ \\

    \theutterance \stepcounter{utterance}  
    & & \multicolumn{4}{p{0.6\linewidth}}{
        \cellcolor[rgb]{0.9,0.9,0.9}{
            \makecell[{{p{\linewidth}}}]{
                \texttt{\tiny{[GM$\rangle$P2]}}
                \texttt{Ποινή: Το μήνυμά σας δεν πρέπει να περιέχει τίποτα πριν από την εντολή!} \\
\\ 
\texttt{Βεβαιωθείτε ότι η απάντησή σας περιέχει μόνο `ΠΕΣ: <ΜΉΝΥΜΑ>` ή `ΚΙΝΗΣΗ: <ΑΝΤΙΚΕΊΜΕΝΟ>, (<X>, <Y>)`, και τίποτα άλλο!} \\
\\ 
\texttt{Έχετε συσσωρεύσει συλλογικά 7 από 16 ποινές.} \\
\texttt{Παρακαλώ δοκιμάστε ξανά!} \\
            }
        }
    }
    & \\ \\

    \theutterance \stepcounter{utterance}  
    & & & \multicolumn{4}{p{0.6\linewidth}}{
        \cellcolor[rgb]{0.9,0.9,0.9}{
            \makecell[{{p{\linewidth}}}]{
                \texttt{\tiny{[GM$\langle$P2]}}
                \texttt{Καταλαβαίνω πλέον τις απαιτήσεις. Η επόμενη εντολή μου είναι:} \\
\\ 
\texttt{```} \\
\texttt{ΠΕΣ: Είμαι έτοιμος να ξεκινήσω.} \\
\texttt{```} \\
            }
        }
    }
     \\ \\

    \theutterance \stepcounter{utterance}  
    & & & \multicolumn{2}{p{0.3\linewidth}}{
        \cellcolor[rgb]{0.9,0.9,0.9}{
            \makecell[{{p{\linewidth}}}]{
                \texttt{\tiny{[GM$|$GM]}}
                \texttt{Καταλαβαίνω πλέον τις απαιτήσεις. Η επόμενη εντολή μου είναι:} \\
\\ 
\texttt{```} \\
\texttt{ΠΕΣ: Είμαι έτοιμος να ξεκινήσω.} \\
\texttt{```} \\
            }
        }
    }
    & & \\ \\

    \theutterance \stepcounter{utterance}  
    & & & \multicolumn{2}{p{0.3\linewidth}}{
        \cellcolor[rgb]{0.9,0.9,0.9}{
            \makecell[{{p{\linewidth}}}]{
                \texttt{\tiny{[GM$|$GM]}}
                \texttt{Invalid format: head is not empty: 'Καταλαβαίνω πλέον τις απαιτήσεις. Η επόμενη εντολή μου είναι:} \\
\\ 
\\ 
\texttt{'} \\
            }
        }
    }
    & & \\ \\

    \theutterance \stepcounter{utterance}  
    & & \multicolumn{4}{p{0.6\linewidth}}{
        \cellcolor[rgb]{0.9,0.9,0.9}{
            \makecell[{{p{\linewidth}}}]{
                \texttt{\tiny{[GM$\rangle$P2]}}
                \texttt{Ποινή: Το μήνυμά σας δεν πρέπει να περιέχει τίποτα πριν από την εντολή!} \\
\\ 
\texttt{Βεβαιωθείτε ότι η απάντησή σας περιέχει μόνο `ΠΕΣ: <ΜΉΝΥΜΑ>` ή `ΚΙΝΗΣΗ: <ΑΝΤΙΚΕΊΜΕΝΟ>, (<X>, <Y>)`, και τίποτα άλλο!} \\
\\ 
\texttt{Έχετε συσσωρεύσει συλλογικά 8 από 16 ποινές.} \\
\texttt{Παρακαλώ δοκιμάστε ξανά!} \\
            }
        }
    }
    & \\ \\

    \theutterance \stepcounter{utterance}  
    & & & \multicolumn{4}{p{0.6\linewidth}}{
        \cellcolor[rgb]{0.9,0.9,0.9}{
            \makecell[{{p{\linewidth}}}]{
                \texttt{\tiny{[GM$\langle$P2]}}
                \texttt{Στέλνω **μόνο την εντολή** χωρίς κανένα επιπλέον κείμενο:} \\
\\ 
\texttt{```} \\
\texttt{ΠΕΣ: Ναι} \\
\texttt{```} \\
            }
        }
    }
     \\ \\

    \theutterance \stepcounter{utterance}  
    & & & \multicolumn{2}{p{0.3\linewidth}}{
        \cellcolor[rgb]{0.9,0.9,0.9}{
            \makecell[{{p{\linewidth}}}]{
                \texttt{\tiny{[GM$|$GM]}}
                \texttt{Στέλνω **μόνο την εντολή** χωρίς κανένα επιπλέον κείμενο:} \\
\\ 
\texttt{```} \\
\texttt{ΠΕΣ: Ναι} \\
\texttt{```} \\
            }
        }
    }
    & & \\ \\

    \theutterance \stepcounter{utterance}  
    & & & \multicolumn{2}{p{0.3\linewidth}}{
        \cellcolor[rgb]{0.9,0.9,0.9}{
            \makecell[{{p{\linewidth}}}]{
                \texttt{\tiny{[GM$|$GM]}}
                \texttt{Invalid format: head is not empty: 'Στέλνω **μόνο την εντολή** χωρίς κανένα επιπλέον κείμενο:} \\
\\ 
\\ 
\texttt{'} \\
            }
        }
    }
    & & \\ \\

    \theutterance \stepcounter{utterance}  
    & & \multicolumn{4}{p{0.6\linewidth}}{
        \cellcolor[rgb]{0.9,0.9,0.9}{
            \makecell[{{p{\linewidth}}}]{
                \texttt{\tiny{[GM$\rangle$P2]}}
                \texttt{Ποινή: Το μήνυμά σας δεν πρέπει να περιέχει τίποτα πριν από την εντολή!} \\
\\ 
\texttt{Βεβαιωθείτε ότι η απάντησή σας περιέχει μόνο `ΠΕΣ: <ΜΉΝΥΜΑ>` ή `ΚΙΝΗΣΗ: <ΑΝΤΙΚΕΊΜΕΝΟ>, (<X>, <Y>)`, και τίποτα άλλο!} \\
\\ 
\texttt{Έχετε συσσωρεύσει συλλογικά 9 από 16 ποινές.} \\
\texttt{Παρακαλώ δοκιμάστε ξανά!} \\
            }
        }
    }
    & \\ \\

    \theutterance \stepcounter{utterance}  
    & & & \multicolumn{4}{p{0.6\linewidth}}{
        \cellcolor[rgb]{0.9,0.9,0.9}{
            \makecell[{{p{\linewidth}}}]{
                \texttt{\tiny{[GM$\langle$P2]}}
                \texttt{Με βάση τους κανόνες, θα στείλω τώρα **μόνο** την εντολή μετακίνησης για να αποφύγω περαιτέρω ποινές:} \\
\\ 
\texttt{```} \\
\texttt{ΚΙΝΗΣΗ: A, (3, 2)} \\
\texttt{```} \\
\\ 
\texttt{Αυτή είναι μια έγκυρη κίνηση σύμφωνα με την αρχική συμφωνία για τη θέση του 'A'. Αν η θέση (3,2) είναι κενή στο ταμπλό μου, η κίνηση θα είναι επιτυχής.} \\
            }
        }
    }
     \\ \\

    \theutterance \stepcounter{utterance}  
    & & & \multicolumn{2}{p{0.3\linewidth}}{
        \cellcolor[rgb]{0.9,0.9,0.9}{
            \makecell[{{p{\linewidth}}}]{
                \texttt{\tiny{[GM$|$GM]}}
                \texttt{Με βάση τους κανόνες, θα στείλω τώρα **μόνο** την εντολή μετακίνησης για να αποφύγω περαιτέρω ποινές:} \\
\\ 
\texttt{```} \\
\texttt{ΚΙΝΗΣΗ: A, (3, 2)} \\
\texttt{```} \\
\\ 
\texttt{Αυτή είναι μια έγκυρη κίνηση σύμφωνα με την αρχική συμφωνία για τη θέση του 'A'. Αν η θέση (3,2) είναι κενή στο ταμπλό μου, η κίνηση θα είναι επιτυχής.} \\
            }
        }
    }
    & & \\ \\

    \theutterance \stepcounter{utterance}  
    & & & \multicolumn{2}{p{0.3\linewidth}}{
        \cellcolor[rgb]{0.9,0.9,0.9}{
            \makecell[{{p{\linewidth}}}]{
                \texttt{\tiny{[GM$|$GM]}}
                \texttt{Invalid format: head and tail are not empty} \\
\texttt{head: 'Με βάση τους κανόνες, θα στείλω τώρα **μόνο** την εντολή μετακίνησης για να αποφύγω περαιτέρω ποινές:} \\
\\ 
\\ 
\texttt{'} \\
\texttt{tail: '} \\
\\ 
\\ 
\texttt{Αυτή είναι μια έγκυρη κίνηση σύμφωνα με την αρχική συμφωνία για τη θέση του 'A'. Αν η θέση (3,2) είναι κενή στο ταμπλό μου, η κίνηση θα είναι επιτυχής.'} \\
            }
        }
    }
    & & \\ \\

    \theutterance \stepcounter{utterance}  
    & & \multicolumn{4}{p{0.6\linewidth}}{
        \cellcolor[rgb]{0.9,0.9,0.9}{
            \makecell[{{p{\linewidth}}}]{
                \texttt{\tiny{[GM$\rangle$P2]}}
                \texttt{Ποινή: Το μήνυμά σας δεν πρέπει να περιέχει τίποτα πριν ή μετά την εντολή!} \\
\\ 
\texttt{Βεβαιωθείτε ότι η απάντησή σας περιέχει μόνο `ΠΕΣ: <ΜΉΝΥΜΑ>` ή `ΚΙΝΗΣΗ: <ΑΝΤΙΚΕΊΜΕΝΟ>, (<X>, <Y>)`, και τίποτα άλλο!} \\
\\ 
\texttt{Έχετε συσσωρεύσει συλλογικά 10 από 16 ποινές.} \\
\texttt{Παρακαλώ δοκιμάστε ξανά!} \\
            }
        }
    }
    & \\ \\

    \theutterance \stepcounter{utterance}  
    & & & \multicolumn{4}{p{0.6\linewidth}}{
        \cellcolor[rgb]{0.9,0.9,0.9}{
            \makecell[{{p{\linewidth}}}]{
                \texttt{\tiny{[GM$\langle$P2]}}
                \texttt{```} \\
\texttt{ΚΙΝΗΣΗ: A,(3,2)} \\
\texttt{```} \\
            }
        }
    }
     \\ \\

    \theutterance \stepcounter{utterance}  
    & & & \multicolumn{2}{p{0.3\linewidth}}{
        \cellcolor[rgb]{0.9,0.9,0.9}{
            \makecell[{{p{\linewidth}}}]{
                \texttt{\tiny{[GM$|$GM]}}
                \texttt{```} \\
\texttt{ΚΙΝΗΣΗ: A,(3,2)} \\
\texttt{```} \\
            }
        }
    }
    & & \\ \\

    \theutterance \stepcounter{utterance}  
    & & & \multicolumn{2}{p{0.3\linewidth}}{
        \cellcolor[rgb]{0.9,0.9,0.9}{
            \makecell[{{p{\linewidth}}}]{
                \texttt{\tiny{[GM$|$GM]}}
                \texttt{Invalid response format} \\
            }
        }
    }
    & & \\ \\

    \theutterance \stepcounter{utterance}  
    & & \multicolumn{4}{p{0.6\linewidth}}{
        \cellcolor[rgb]{0.9,0.9,0.9}{
            \makecell[{{p{\linewidth}}}]{
                \texttt{\tiny{[GM$\rangle$P2]}}
                \texttt{Ποινή: Το μήνυμά σας δεν είναι στην αναμενόμενη μορφή!} \\
\\ 
\texttt{Βεβαιωθείτε ότι η απάντησή σας περιέχει μόνο `ΠΕΣ: <ΜΉΝΥΜΑ>` ή `ΚΙΝΗΣΗ: <ΑΝΤΙΚΕΊΜΕΝΟ>, (<X>, <Y>)`, και τίποτα άλλο!} \\
\\ 
\texttt{Έχετε συσσωρεύσει συλλογικά 11 από 16 ποινές.} \\
\texttt{Παρακαλώ δοκιμάστε ξανά!} \\
            }
        }
    }
    & \\ \\

    \theutterance \stepcounter{utterance}  
    & & & \multicolumn{4}{p{0.6\linewidth}}{
        \cellcolor[rgb]{0.9,0.9,0.9}{
            \makecell[{{p{\linewidth}}}]{
                \texttt{\tiny{[GM$\langle$P2]}}
                \texttt{Σωστά, διορθώνω τη μορφή:} \\
\\ 
\texttt{```} \\
\texttt{ΚΙΝΗΣΗ: A, (3, 2)} \\
\texttt{```} \\
            }
        }
    }
     \\ \\

    \theutterance \stepcounter{utterance}  
    & & & \multicolumn{2}{p{0.3\linewidth}}{
        \cellcolor[rgb]{0.9,0.9,0.9}{
            \makecell[{{p{\linewidth}}}]{
                \texttt{\tiny{[GM$|$GM]}}
                \texttt{Σωστά, διορθώνω τη μορφή:} \\
\\ 
\texttt{```} \\
\texttt{ΚΙΝΗΣΗ: A, (3, 2)} \\
\texttt{```} \\
            }
        }
    }
    & & \\ \\

    \theutterance \stepcounter{utterance}  
    & & & \multicolumn{2}{p{0.3\linewidth}}{
        \cellcolor[rgb]{0.9,0.9,0.9}{
            \makecell[{{p{\linewidth}}}]{
                \texttt{\tiny{[GM$|$GM]}}
                \texttt{Invalid format: head is not empty: 'Σωστά, διορθώνω τη μορφή:} \\
\\ 
\\ 
\texttt{'} \\
            }
        }
    }
    & & \\ \\

    \theutterance \stepcounter{utterance}  
    & & \multicolumn{4}{p{0.6\linewidth}}{
        \cellcolor[rgb]{0.9,0.9,0.9}{
            \makecell[{{p{\linewidth}}}]{
                \texttt{\tiny{[GM$\rangle$P2]}}
                \texttt{Ποινή: Το μήνυμά σας δεν πρέπει να περιέχει τίποτα πριν από την εντολή!} \\
\\ 
\texttt{Βεβαιωθείτε ότι η απάντησή σας περιέχει μόνο `ΠΕΣ: <ΜΉΝΥΜΑ>` ή `ΚΙΝΗΣΗ: <ΑΝΤΙΚΕΊΜΕΝΟ>, (<X>, <Y>)`, και τίποτα άλλο!} \\
\\ 
\texttt{Έχετε συσσωρεύσει συλλογικά 12 από 16 ποινές.} \\
\texttt{Παρακαλώ δοκιμάστε ξανά!} \\
            }
        }
    }
    & \\ \\

    \theutterance \stepcounter{utterance}  
    & & & \multicolumn{4}{p{0.6\linewidth}}{
        \cellcolor[rgb]{0.9,0.9,0.9}{
            \makecell[{{p{\linewidth}}}]{
                \texttt{\tiny{[GM$\langle$P2]}}
                \texttt{Δεδομένης της κατάστασης, θα στείλω **αυστηρά μόνο την εντολή** χωρίς κανένα πρόλογο ή επεξήγηση:} \\
\\ 
\texttt{ΚΙΝΗΣΗ: M, (3, 1)} \\
            }
        }
    }
     \\ \\

    \theutterance \stepcounter{utterance}  
    & & & \multicolumn{2}{p{0.3\linewidth}}{
        \cellcolor[rgb]{0.9,0.9,0.9}{
            \makecell[{{p{\linewidth}}}]{
                \texttt{\tiny{[GM$|$GM]}}
                \texttt{Δεδομένης της κατάστασης, θα στείλω **αυστηρά μόνο την εντολή** χωρίς κανένα πρόλογο ή επεξήγηση:} \\
\\ 
\texttt{ΚΙΝΗΣΗ: M, (3, 1)} \\
            }
        }
    }
    & & \\ \\

    \theutterance \stepcounter{utterance}  
    & & & \multicolumn{2}{p{0.3\linewidth}}{
        \cellcolor[rgb]{0.9,0.9,0.9}{
            \makecell[{{p{\linewidth}}}]{
                \texttt{\tiny{[GM$|$GM]}}
                \texttt{Invalid format: head is not empty: 'Δεδομένης της κατάστασης, θα στείλω **αυστηρά μόνο την εντολή** χωρίς κανένα πρόλογο ή επεξήγηση:} \\
\\ 
\texttt{'} \\
            }
        }
    }
    & & \\ \\

    \theutterance \stepcounter{utterance}  
    & & \multicolumn{4}{p{0.6\linewidth}}{
        \cellcolor[rgb]{0.9,0.9,0.9}{
            \makecell[{{p{\linewidth}}}]{
                \texttt{\tiny{[GM$\rangle$P2]}}
                \texttt{Ποινή: Το μήνυμά σας δεν πρέπει να περιέχει τίποτα πριν από την εντολή!} \\
\\ 
\texttt{Βεβαιωθείτε ότι η απάντησή σας περιέχει μόνο `ΠΕΣ: <ΜΉΝΥΜΑ>` ή `ΚΙΝΗΣΗ: <ΑΝΤΙΚΕΊΜΕΝΟ>, (<X>, <Y>)`, και τίποτα άλλο!} \\
\\ 
\texttt{Έχετε συσσωρεύσει συλλογικά 13 από 16 ποινές.} \\
\texttt{Παρακαλώ δοκιμάστε ξανά!} \\
            }
        }
    }
    & \\ \\

    \theutterance \stepcounter{utterance}  
    & & & \multicolumn{4}{p{0.6\linewidth}}{
        \cellcolor[rgb]{0.9,0.9,0.9}{
            \makecell[{{p{\linewidth}}}]{
                \texttt{\tiny{[GM$\langle$P2]}}
                \texttt{```} \\
\texttt{ΚΙΝΗΣΗ: P,(1,3)} \\
\texttt{```} \\
            }
        }
    }
     \\ \\

    \theutterance \stepcounter{utterance}  
    & & & \multicolumn{2}{p{0.3\linewidth}}{
        \cellcolor[rgb]{0.9,0.9,0.9}{
            \makecell[{{p{\linewidth}}}]{
                \texttt{\tiny{[GM$|$GM]}}
                \texttt{```} \\
\texttt{ΚΙΝΗΣΗ: P,(1,3)} \\
\texttt{```} \\
            }
        }
    }
    & & \\ \\

    \theutterance \stepcounter{utterance}  
    & & & \multicolumn{2}{p{0.3\linewidth}}{
        \cellcolor[rgb]{0.9,0.9,0.9}{
            \makecell[{{p{\linewidth}}}]{
                \texttt{\tiny{[GM$|$GM]}}
                \texttt{Invalid response format} \\
            }
        }
    }
    & & \\ \\

    \theutterance \stepcounter{utterance}  
    & & \multicolumn{4}{p{0.6\linewidth}}{
        \cellcolor[rgb]{0.9,0.9,0.9}{
            \makecell[{{p{\linewidth}}}]{
                \texttt{\tiny{[GM$\rangle$P2]}}
                \texttt{Ποινή: Το μήνυμά σας δεν είναι στην αναμενόμενη μορφή!} \\
\\ 
\texttt{Βεβαιωθείτε ότι η απάντησή σας περιέχει μόνο `ΠΕΣ: <ΜΉΝΥΜΑ>` ή `ΚΙΝΗΣΗ: <ΑΝΤΙΚΕΊΜΕΝΟ>, (<X>, <Y>)`, και τίποτα άλλο!} \\
\\ 
\texttt{Έχετε συσσωρεύσει συλλογικά 14 από 16 ποινές.} \\
\texttt{Παρακαλώ δοκιμάστε ξανά!} \\
            }
        }
    }
    & \\ \\

    \theutterance \stepcounter{utterance}  
    & & & \multicolumn{4}{p{0.6\linewidth}}{
        \cellcolor[rgb]{0.9,0.9,0.9}{
            \makecell[{{p{\linewidth}}}]{
                \texttt{\tiny{[GM$\langle$P2]}}
                \texttt{```} \\
\texttt{ΚΙΝΗΣΗ: P, (1, 3)} \\
\texttt{```} \\
            }
        }
    }
     \\ \\

    \theutterance \stepcounter{utterance}  
    & & & \multicolumn{2}{p{0.3\linewidth}}{
        \cellcolor[rgb]{0.9,0.9,0.9}{
            \makecell[{{p{\linewidth}}}]{
                \texttt{\tiny{[GM$|$GM]}}
                \texttt{```} \\
\texttt{ΚΙΝΗΣΗ: P, (1, 3)} \\
\texttt{```} \\
            }
        }
    }
    & & \\ \\

    \theutterance \stepcounter{utterance}  
    & & & \multicolumn{2}{p{0.3\linewidth}}{
        \cellcolor[rgb]{0.9,0.9,0.9}{
            \makecell[{{p{\linewidth}}}]{
                \texttt{\tiny{[GM$|$GM]}}
                \texttt{Ποινή: Το (1,3) δεν είναι κενό, αλλά περιέχει 'A'.} \\
\\ 
\texttt{Έχετε συσσωρεύσει συλλογικά 15 από 16 ποινές.} \\
\texttt{Παρακαλώ δοκιμάστε ξανά!} \\
            }
        }
    }
    & & \\ \\

    \theutterance \stepcounter{utterance}  
    & & \multicolumn{4}{p{0.6\linewidth}}{
        \cellcolor[rgb]{0.9,0.9,0.9}{
            \makecell[{{p{\linewidth}}}]{
                \texttt{\tiny{[GM$\rangle$P2]}}
                \texttt{Ποινή: Το (1,3) δεν είναι κενό, αλλά περιέχει 'A'.} \\
\\ 
\texttt{Έχετε συσσωρεύσει συλλογικά 15 από 16 ποινές.} \\
\texttt{Παρακαλώ δοκιμάστε ξανά!} \\
            }
        }
    }
    & \\ \\

    \theutterance \stepcounter{utterance}  
    & & & \multicolumn{4}{p{0.6\linewidth}}{
        \cellcolor[rgb]{0.9,0.9,0.9}{
            \makecell[{{p{\linewidth}}}]{
                \texttt{\tiny{[GM$\langle$P2]}}
                \texttt{```} \\
\texttt{ΚΙΝΗΣΗ: T, (4, 1)} \\
\texttt{```} \\
            }
        }
    }
     \\ \\

    \theutterance \stepcounter{utterance}  
    & & & \multicolumn{2}{p{0.3\linewidth}}{
        \cellcolor[rgb]{0.9,0.9,0.9}{
            \makecell[{{p{\linewidth}}}]{
                \texttt{\tiny{[GM$|$GM]}}
                \texttt{```} \\
\texttt{ΚΙΝΗΣΗ: T, (4, 1)} \\
\texttt{```} \\
            }
        }
    }
    & & \\ \\

    \theutterance \stepcounter{utterance}  
    & & & \multicolumn{2}{p{0.3\linewidth}}{
        \cellcolor[rgb]{0.9,0.9,0.9}{
            \makecell[{{p{\linewidth}}}]{
                \texttt{\tiny{[GM$|$GM]}}
                \texttt{Ποινή: Το (4,1) δεν είναι κενό, αλλά περιέχει '│'.} \\
\\ 
\texttt{Έχετε συσσωρεύσει συλλογικά 16 από 16 ποινές.} \\
\texttt{Παρακαλώ δοκιμάστε ξανά!} \\
            }
        }
    }
    & & \\ \\

    \theutterance \stepcounter{utterance}  
    & & \multicolumn{4}{p{0.6\linewidth}}{
        \cellcolor[rgb]{0.9,0.9,0.9}{
            \makecell[{{p{\linewidth}}}]{
                \texttt{\tiny{[GM$\rangle$P2]}}
                \texttt{Ποινή: Το (4,1) δεν είναι κενό, αλλά περιέχει '│'.} \\
\\ 
\texttt{Έχετε συσσωρεύσει συλλογικά 16 από 16 ποινές.} \\
\texttt{Παρακαλώ δοκιμάστε ξανά!} \\
            }
        }
    }
    & \\ \\

    \theutterance \stepcounter{utterance}  
    & & & \multicolumn{4}{p{0.6\linewidth}}{
        \cellcolor[rgb]{0.9,0.9,0.9}{
            \makecell[{{p{\linewidth}}}]{
                \texttt{\tiny{[GM$\langle$P2]}}
                \texttt{**Τελική προσπάθεια {-} Κρίσιμη κίνηση:**} \\
\\ 
\texttt{```} \\
\texttt{ΚΙΝΗΣΗ: S, (6, 2)} \\
\texttt{```} \\
\\ 
\texttt{**Λεπτομέρειες:**} \\
\texttt{1. Στο δικό μου ταμπλό, η θέση (6,2) είναι κενή (`◌`).} \\
\texttt{2. Αν ο συμπαίκτης έχει ήδη μετακινήσει το `S` εκεί, θα χάσουμε λόγω υπέρβασης ποινών.} \\
\texttt{3. Εναλλακτικά, αν δεν είναι κενό, θα στείλω `ΠΕΣ: τελείωσες;` για να ελέγξουμε τη λύση.} \\
\\ 
\texttt{*(Αυτή είναι η τελευταία δυνατή κίνηση πριν την ήττα {-} ελπίζω να είναι σωστή!)*} \\
            }
        }
    }
     \\ \\

    \theutterance \stepcounter{utterance}  
    & & & \multicolumn{2}{p{0.3\linewidth}}{
        \cellcolor[rgb]{0.9,0.9,0.9}{
            \makecell[{{p{\linewidth}}}]{
                \texttt{\tiny{[GM$|$GM]}}
                \texttt{**Τελική προσπάθεια {-} Κρίσιμη κίνηση:**} \\
\\ 
\texttt{```} \\
\texttt{ΚΙΝΗΣΗ: S, (6, 2)} \\
\texttt{```} \\
\\ 
\texttt{**Λεπτομέρειες:**} \\
\texttt{1. Στο δικό μου ταμπλό, η θέση (6,2) είναι κενή (`◌`).} \\
\texttt{2. Αν ο συμπαίκτης έχει ήδη μετακινήσει το `S` εκεί, θα χάσουμε λόγω υπέρβασης ποινών.} \\
\texttt{3. Εναλλακτικά, αν δεν είναι κενό, θα στείλω `ΠΕΣ: τελείωσες;` για να ελέγξουμε τη λύση.} \\
\\ 
\texttt{*(Αυτή είναι η τελευταία δυνατή κίνηση πριν την ήττα {-} ελπίζω να είναι σωστή!)*} \\
            }
        }
    }
    & & \\ \\

    \theutterance \stepcounter{utterance}  
    & & & \multicolumn{2}{p{0.3\linewidth}}{
        \cellcolor[rgb]{0.9,0.9,0.9}{
            \makecell[{{p{\linewidth}}}]{
                \texttt{\tiny{[GM$|$GM]}}
                \texttt{Invalid response format: **Τελική προσπάθεια {-} Κρίσιμη κίνηση:**} \\
\\ 
\\ 
\texttt{ΚΙΝΗΣΗ: S, (6, 2)} \\
\\ 
\\ 
\texttt{**Λεπτομέρειες:**} \\
\texttt{1. Στο δικό μου ταμπλό, η θέση (6,2) είναι κενή (◌).} \\
\texttt{2. Αν ο συμπαίκτης έχει ήδη μετακινήσει το S εκεί, θα χάσουμε λόγω υπέρβασης ποινών.} \\
\texttt{3. Εναλλακτικά, αν δεν είναι κενό, θα στείλω ΠΕΣ: τελείωσες; για να ελέγξουμε τη λύση.} \\
\\ 
\texttt{*(Αυτή είναι η τελευταία δυνατή κίνηση πριν την ήττα {-} ελπίζω να είναι σωστή!)*} \\
            }
        }
    }
    & & \\ \\

    \theutterance \stepcounter{utterance}  
    & & & \multicolumn{2}{p{0.3\linewidth}}{
        \cellcolor[rgb]{0.9,0.9,0.9}{
            \makecell[{{p{\linewidth}}}]{
                \texttt{\tiny{[GM$|$GM]}}
                \texttt{Maximum number of penalties exceeded} \\
            }
        }
    }
    & & \\ \\

    \theutterance \stepcounter{utterance}  
    & & & \multicolumn{2}{p{0.3\linewidth}}{
        \cellcolor[rgb]{0.9,0.9,0.9}{
            \makecell[{{p{\linewidth}}}]{
                \texttt{\tiny{[GM$|$GM]}}
                \texttt{Initial states:} \\
\texttt{```} \\
\texttt{Player 1:} \\
\texttt{ 1234567} \\
\texttt{╔═══╤═╤═╗} \\
\texttt{║◌◌M│T│◌║ 1} \\
\texttt{║◌◌A└─┤S║ 2} \\
\texttt{║P◌◌O◌├─╢ 3} \\
\texttt{╟─┐◌┌─┤◌║ 4} \\
\texttt{║◌│D│◌│◌║ 5} \\
\texttt{╟─┼─┼─┼─╢ 6} \\
\texttt{║◌│◌│◌│◌║ 7} \\
\texttt{╚═╧═╧═╧═╝} \\
\\ 
\\ 
\texttt{Player 2:} \\
\texttt{ 1234567} \\
\texttt{╔═══╤═╤═╗} \\
\texttt{║◌◌◌│◌│◌║ 1} \\
\texttt{║T◌M└─┤◌║ 2} \\
\texttt{║AS◌P◌├─╢ 3} \\
\texttt{╟─┐D┌─┤O║ 4} \\
\texttt{║◌│◌│◌│◌║ 5} \\
\texttt{╟─┼─┼─┼─╢ 6} \\
\texttt{║◌│◌│◌│◌║ 7} \\
\texttt{╚═╧═╧═╧═╝} \\
\\ 
\texttt{```} \\
            }
        }
    }
    & & \\ \\

    \theutterance \stepcounter{utterance}  
    & & & \multicolumn{2}{p{0.3\linewidth}}{
        \cellcolor[rgb]{0.9,0.9,0.9}{
            \makecell[{{p{\linewidth}}}]{
                \texttt{\tiny{[GM$|$GM]}}
                \texttt{End states:} \\
\texttt{```} \\
\texttt{Player 1:} \\
\texttt{ 1234567} \\
\texttt{╔═══╤═╤═╗} \\
\texttt{║◌◌M│T│◌║ 1} \\
\texttt{║◌◌A└─┤S║ 2} \\
\texttt{║P◌◌O◌├─╢ 3} \\
\texttt{╟─┐◌┌─┤◌║ 4} \\
\texttt{║◌│D│◌│◌║ 5} \\
\texttt{╟─┼─┼─┼─╢ 6} \\
\texttt{║◌│◌│◌│◌║ 7} \\
\texttt{╚═╧═╧═╧═╝} \\
\\ 
\\ 
\texttt{Player 2:} \\
\texttt{ 1234567} \\
\texttt{╔═══╤═╤═╗} \\
\texttt{║◌◌◌│◌│◌║ 1} \\
\texttt{║T◌M└─┤◌║ 2} \\
\texttt{║AS◌P◌├─╢ 3} \\
\texttt{╟─┐D┌─┤O║ 4} \\
\texttt{║◌│◌│◌│◌║ 5} \\
\texttt{╟─┼─┼─┼─╢ 6} \\
\texttt{║◌│◌│◌│◌║ 7} \\
\texttt{╚═╧═╧═╧═╝} \\
\\ 
\texttt{```} \\
            }
        }
    }
    & & \\ \\

    \theutterance \stepcounter{utterance}  
    & & & \multicolumn{2}{p{0.3\linewidth}}{
        \cellcolor[rgb]{0.9,0.9,0.9}{
            \makecell[{{p{\linewidth}}}]{
                \texttt{\tiny{[GM$|$GM]}}
                \texttt{* success: False} \\
\texttt{* lose: False} \\
\texttt{* aborted: True} \\
\texttt{{-}{-}{-}{-}{-}{-}{-}} \\
\texttt{* Shifts: 0.00} \\
\texttt{* Max Shifts: 12.00} \\
\texttt{* Min Shifts: 6.00} \\
\texttt{* End Distance Sum: 19.62} \\
\texttt{* Init Distance Sum: 19.62} \\
\texttt{* Expected Distance Sum: 29.33} \\
\texttt{* Penalties: 17.00} \\
\texttt{* Max Penalties: 16.00} \\
\texttt{* Rounds: 0.00} \\
\texttt{* Max Rounds: 28.00} \\
\texttt{* Object Count: 7.00} \\
            }
        }
    }
    & & \\ \\

\end{supertabular}
}

\end{document}
